\section{Parameter provenance}
\label{sec:provenance}

A valuation model is only as interpretable as the audit trail behind its
inputs. We therefore classify every parameter into one of four categories and
report the classification alongside the value itself.

\begin{description}
\item[Calibrated] determined by observed market prices through the constrained
  optimisation of Section~\ref{sec:curve}.
\item[Inherited] taken from a historical estimation exercise under $\PP$ and
  used for valuation without a measure adjustment.
\item[Derived] reconstructed from summary statistics of that exercise because
  the primary estimate was unavailable, with no standard error attached.
\item[Assumed] chosen by the analyst on economic or pragmatic grounds.
\end{description}

Table~\ref{tab:provenance} applies the scheme. Only the forward curve is
calibrated. Everything governing the distribution around it is inherited,
derived or assumed, which is the concrete form the incompleteness of this
market takes.

\begin{table}[htbp]
\centering
\caption{Production parameters and provenance. Only $F(\cdot)$ is identified by
market prices.}
\label{tab:provenance}
\small
\setlength{\tabcolsep}{4pt}
\begin{tabular}{@{}llll@{}}
\toprule
Parameter & Value & Status & Note \\
\midrule
$F(\cdot)$ & curve & Calibrated & Exact on six monthly averages \\
Jan.\ anchor & \texttt{spot\_to\_next} & Assumed & No quote exists \\
$\kappa$ (h$^{-1}$) & $0.078394$ & Inherited & Half-life $8.84$\,h \\
$\sigma^{y}_{0}$ & $0.0035348$ & Inherited & Calm regime \\
$\sigma^{y}_{1}$ & $0.0924067$ & Inherited & Stress regime \\
$s_P$ (\TRY) & $282.48$ & Assumed & Source file unavailable \\
$\alpha_{01}$ & $-1.015666$ & Derived & Duration reconstruction \\
$\alpha_{10}$ & $-1.895229$ & Derived & Duration reconstruction \\
$\gamma_{01}$ & $-0.583778$ & Inherited & Ramp term omitted \\
$\gamma_{10}$ & $+0.077698$ & Inherited & Ramp term omitted \\
$\pi_{1,0}$ & $0.068023$ & Inherited & Filtered at valuation \\
$a_0,a_1$ & $0,\ 0$ & Assumed & Zero drift premium \\
$r$ (annual) & $0.40$ & Assumed & Flat, not calibrated \\
\bottomrule
\end{tabular}
\end{table}

Three entries require comment because they carry more interpretive weight than
their appearance in a table suggests.

The intercepts $\alpha_{01},\alpha_{10}$ were not available from the original
estimation output, which reported slope coefficients and average marginal
effects but not intercepts. We reconstructed them by solving, over the full
covariate sample, the two scalar equations that match the implied average
departure rates to the reported mean regime durations of $3.76$ and $7.56$
hours. The reconstruction reproduces the reported stationary stress share to
within $1\%$ under an ergodic approximation and $4.5\%$ under simulation along
the realised covariate path. It is nevertheless a moment-matching
reconstruction and not a likelihood estimate, so no standard error attaches to
it and no inference should be drawn from its magnitude.

The slopes $\gamma_{01},\gamma_{10}$ come from a specification that also
contained a residual-demand ramp covariate, which the production implementation
does not carry. If the original fit was joint, our slopes are subject to omitted
variable bias of unknown sign. We were unable to determine from the available
artefacts whether the two covariates were estimated jointly or separately, and
we report the ambiguity rather than resolving it by assumption.

The initial regime distribution $\pi$ comes from a filter run under an earlier
model specification than the one supplying the volatilities, so the production
configuration mixes sources at the initialisation step. We quantify the
consequence directly in Section~\ref{sec:sensitivity} by re-pricing under the
alternative of initialising at the stationary distribution.
